\section{Linguistic Distance Analysis}

\subsection{Motivation}

If Norwegian is structurally closer to Danish in syntax and lexicon, it may offer a stronger transfer signal. The hypothesis sounds reasonable. The subsections below put a number on that closeness, then compare it to the observed F1 differences to see how well the two align.

\subsection{Method}

We compute typological feature vectors with \texttt{lang2vec} \citep{littell2017uriel}, using the \texttt{syntax\_wals} feature set from WALS. This representation encodes approximately 103 syntactic and morphological features; features unavailable for a given language pair are excluded before computing cosine distance. Formally,
\[
d(L_1, L_2) = 1 - \frac{\mathbf{v}_1 \cdot \mathbf{v}_2}{\|\mathbf{v}_1\|\,\|\mathbf{v}_2\|},
\]
where each vector $\mathbf{v}_i$ contains only the jointly available features for that pair.

\subsection{Results}

Table~\ref{tab:language_distance} lists pairwise cosine distances for English, Danish, Norwegian (Bokm\aa l), Swedish and German. Swedish and German serve as reference points: Swedish sits close to Danish within the North Germanic cluster, while German occupies a more distant West Germanic position.

\begin{table}[t]
\centering
\small
\resizebox{\linewidth}{!}{%
\begin{tabular}{lccccc}
\toprule
          & English & Danish & Norwegian & Swedish & German \\
\midrule
English   & 0.000   & 0.119  & 0.109     & 0.080   & 0.145  \\
Danish    & 0.119   & 0.000  & \textbf{0.021}     & 0.039   & 0.124  \\
Norwegian & 0.109   & 0.021  & 0.000     & 0.065   & 0.102  \\
Swedish   & 0.080   & 0.039  & 0.065     & 0.000   & 0.164  \\
German    & 0.145   & 0.124  & 0.102     & 0.164   & 0.000  \\
\bottomrule
\end{tabular}%
}
\caption{Pairwise cosine distances based on WALS syntactic features (lang2vec \texttt{syntax\_wals}). Lower = more similar. Norwegian is the closest language to Danish (0.021), followed by Swedish (0.039); English and German are substantially more distant ($\approx$0.12).}
\label{tab:language_distance}
\end{table}

The ordering is clear:
\begin{itemize}
  \item Norwegian is the closest language to Danish in this space (0.021) — roughly 5.7$\times$ closer than English (0.119).
  \item Swedish is the next closest (0.039), confirming the coherence of the North Germanic subgroup.
  \item English and German fall at comparable distances (0.119 and 0.124), placing both in a distinctly more distant tier relative to Danish.
\end{itemize}

\subsection{Interpretation}

WALS-based cosine distance is a clean, reproducible proxy for structural similarity, but it does not capture the full picture. Pretraining data volume and tokenizer overlap also shape how well a multilingual encoder represents a given language — and those factors are not encoded in WALS vectors. What the distances do support is an unambiguous ordering: NO $<$ SV $<$ EN $\approx$ DE. That ordering is directionally consistent with the NO$\rightarrow$DA $>$ EN$\rightarrow$DA pattern observed in zero-shot experiments, even when individual transfer gaps fall short of statistical significance.
